\scanpage{72}
\mspage{21}

% slanted
\margin{$\times$ on utilise \textcircled{4} sous forme plus forte que $W(4\,\mathrm b'')$ !} $f\in W$, il s'ensuit $f_i\in W$ pour tt $i\in I$.

Si $\pi_0(f)$ n'était pas surjective, prenant $i$ qui n'est pas en l'image, on aurait $X_i=\varnothing\longrightarrow Y_i\ne\varnothing$, qui serait dans $W$. Soit donc \[f:X=\varnothing\longrightarrow Y\] dans $W$. Alors $\forall Z$ dans $\mathrm{Cat}$, $f\times\mathrm{id}_Z:\underset{\substack{\Vert\\\varnothing}}{\underline{X\times Z}}\longrightarrow Y\times Z$ est dans $W$ (car $Z$ est $W$-fibrée sur $e$) donc $\varnothing\longrightarrow Y\times Z$ est dans $W$. Donc aussi la projection \diagram{\begin{tikzcd}[ampersand replacement=\&,column sep=small,row sep=small] Y\times Z\arrow[rr]\&\&Y,\\ \&\varnothing\arrow[ul]\arrow[ur]\&\end{tikzcd}} Comme cette projection est une $W$-fibration, il s'ensuit \deleted{que} par $W(4\,\mathrm b'')$ que ses fibres sont $W$-asphériques. Comme $Y\ne\varnothing$, il y a une telle fibre, elle est $\simeq Z$. Donc $Z$ est $W$-asphérique. \deleted{Soit $f:X$} Donc on trouve que \uline{toute} catégorie \inserted{$Z$ de $\mathrm{Cat}$} est $W$-asphérique. Donc toute flèche $Z\longrightarrow Z'$ de $\mathrm{Cat}$ est dans $W$, \unsure{cqed|cqfd}.

Supposons maintenant $\pi_0(f)$ surjective, \deleted{mais} non injective. Il y aurait une fibre de $\pi_0(X)\longrightarrow\pi_0(Y)=I$ qui serait de cardinal $>1$. En prenant le $f_i$ correspondant, on trouve $X_i\longrightarrow Y_i$ avec $Y_i$ 0-connexe, et $\mathrm{card}\,\pi_0(X_i)\geq2$. \deleted{\unsure{et}} Donc on est ramené au cas \[f:X\longrightarrow Y\qquad\left\{\begin{aligned}&f\in W\\&Y\text{ 0-connexe}\\&\mathrm{Card}\,\pi_0(X)\geq2.\end{aligned}\right.\]

Considérons la factorisation canonique \[X\xrightarrow{i_f}\overline X\xrightarrow{\overline f}Y,\] avec $i_f\in W$, donc $\overline f\in W$. Donc par $(W\,4\,\mathrm b'')$ les fibres de $\overline f$ sont $W$-asphériques. Mais on voit tout de suite, \deleted{que} en décomposant $\overline X$ suivant les comp. conn. $X_i$ de $X$, que ces fibres ont un $\pi_0$ de card $\geq\mathrm{Card}\,\pi_0(X)\geq2$. Donc on est ramené au cas où $Y=e$, i.e \[X\quad W\text{-asphérique,}\qquad \mathrm{card}\,\pi_0(X)\geq2.\] Considérons $\unsure{X}\times X$, qui est aussi $W$ asphérique, et l'application diagonale

% archive sheet number: 72
